\documentclass[compress, xcolor={svgnames}, aspectratio=169, smaller]{beamer}
\newcommand{\CommonPath}{../../Common}
\newcommand{\coursename}{Intro to DL - S2}
\usepackage{\CommonPath/slidespreamble}
\usepackage{\CommonPath/math_commands}
% \setbeameroption{show notes}
\setbeameroption{hide notes}

\title{Introduction to Deep Learning}
\subtitle{Session 2: Perceptron, MLPs \& Activations \\[1em]
\small March 2026}
\author{Magda Gregorová - \href{mailto:magda.gregorova@thws.de}{magda.gregorova@thws.de}}
\institute{
\includegraphics[width=0.3\textwidth]{thws-logo_vert_en_orange-rgb.png}}
\hypersetup{
pdftitle={Introduction to Deep Learning},
pdfauthor={Magda Gregorová},
pdfkeywords={perceptron, MLP, activation functions, backpropagation}
}

\date{{\color{skyblue1}\footnotesize Licensed according to
\href{https://creativecommons.org/licenses/by-sa/4.0}{\color{skyblue1}CC-BY-SA 4.0}}}

\begin{document}

\begin{frame}[plain]
    \titlepage
\end{frame}

\begin{frame}{Supervised learning}
\begin{columns}
    \begin{column}{0.48\textwidth}
        \structure{Classification (binary)}
        \begin{itemize}
            \item Output: discrete label $y \in \{0, 1\}$
            \item Goal: find $f: \R^d \to \{0, 1\}$
        \end{itemize}
        \vskip 0.5em
        \centering
        % Common/tikz/classification_2d.tex
\begin{tikzpicture}[scale=0.9]
\tikzstyle{posit} = [circle, draw, fill=thwsorange!60, minimum size=0.1cm, inner sep=0pt]
\tikzstyle{neg} = [circle, draw, fill=thwsblue!40, minimum size=0.1cm, inner sep=0pt]
\tikzstyle{annot} = [font=\scriptsize]

% Axes
\draw[-Stealth, thick] (-0.3, 0) -- (3.8, 0) node[right, annot] {$x_1$};
\draw[-Stealth, thick] (0, -0.3) -- (0, 3.2) node[above, annot] {$x_2$};

% Classificaiton boundary
\onslide<2->{
\draw[thick, thwsorange] (0.0, 0.3) -- (3.6, 2.6) node[annot, anchor=south] {separating line};}

% posititive class
\node[posit] at (0.5, 2.4) {};
\node[posit] at (0.8, 2.9) {};
\node[posit] at (0.4, 2.0) {};
\node[posit] at (1.1, 2.6) {};
\node[posit] at (0.7, 3.0) {};

% Negative class
\node[neg] at (2.6, 0.5) {};
\node[neg] at (3.0, 0.9) {};
\node[neg] at (3.2, 0.4) {};
\node[neg] at (2.4, 1.0) {};
\node[neg] at (2.9, 1.3) {};

% Labels
\node[annot, thwsorange] at (1.7, 2.9) {$y=1$};
\node[annot, thwsblue]   at (1.8, 0.5) {$y=0$};


\end{tikzpicture}
    \end{column}
    \begin{column}{0.48\textwidth}
        \structure{Regression}
        \begin{itemize}
            \item Output: continuous value $y \in \R$
            \item Goal: find $f: \R^d \to \R$
        \end{itemize}
        \vskip 0.5em
        \centering
        % Common/tikz/regression_1d.tex
\begin{tikzpicture}[scale=0.9]
\tikzstyle{pt}    = [circle, draw, fill=thwsblue!40, minimum size=0.1cm, inner sep=0pt]
\tikzstyle{annot} = [font=\scriptsize]

% Axes
\draw[-Stealth, thick] (-0.3, 0) -- (3.8, 0) node[right, annot] {$x$};
\draw[-Stealth, thick] (0, -0.3) -- (0, 3.2) node[above, annot] {$y$};

% % Regression line
% \onslide<2->{
% \draw[thick, thwsorange] (0.1, 0.4) -- (3.6, 2.6) node[annot, right, anchor=south] {regression line};}

% Data points scattered around curve
\node[pt] at (0.3, 0.6) {};
\node[pt] at (0.7, 0.9) {};
\node[pt] at (1.0, 0.8) {};
\node[pt] at (1.3, 1.2) {};
\node[pt] at (1.6, 1.5) {};
\node[pt] at (1.9, 1.4) {};
\node[pt] (yx) at (2.2, 1.8) {};
\node[pt] at (2.5, 2.2) {};
\node[pt] at (2.8, 2.1) {};
\node[pt] at (3.2, 2.2) {};
\node[pt] at (3.5, 2.5) {};

% y point
\draw[dotted, thick] (yx) -- ++(0, -1.9) node[below, annot] {$x^{(i)}$};
\draw[dotted, thick] (yx) -- ++(-2.2, 0) node[left, annot] {$y^{(i)}$};;

% Label
% \node[annot, thwsorange] at (2.8, 2.8) {$f(x)$};
\end{tikzpicture}
    \end{column}
\end{columns}
\end{frame}

\begin{frame}{Linear classifier}
\begin{columns}
    \begin{column}{0.55\textwidth}
        \begin{itemize}
            \item Line in 2D: 
            \begin{align*}
                x_2 & = m x_1 + c \\
                0 & = m x_1 + c - x_2
            \end{align*}
            \item<2-> Separating line: 
            \begin{align*}
                w_1 x_1 + w_2 x_2 + b & = 0\\
                \vw\T\vx + b & = 0
            \end{align*}
            \item<3-> Decision rule:
            \[ \hat{y} = \begin{cases} 1 & \vw\T\vx + b \geq 0 \\ 0 & \text{otherwise} \end{cases} \]
        \end{itemize}
    \end{column}
    \begin{column}{0.45\textwidth}
        \centering
        \input{\CommonPath/tikz/classification_2d_math.tex}
    \end{column}
\end{columns}
\end{frame}

\begin{frame}{From biological to artificial neuron}
\begin{columns}[t]
    \begin{column}{0.48\textwidth}
        \structure{Biological neuron}
        \begin{itemize}
            \item Dendrites receive signals
            \item Cell body integrates inputs
            \item Axon fires if threshold exceeded
        \end{itemize}
        \vskip 1em
\centering
        \includegraphics[width=0.9\textwidth]{Neuron3}\\[-0.5em]
        {\tiny Image: Egm4313.s12 at English Wikipedia, \href{https://creativecommons.org/licenses/by-sa/3.0}{CC BY-SA 3.0}, via Wikimedia Commons}
    \end{column}
    \begin{column}{0.5\textwidth}
        \structure{Artificial neuron $f: \R^d \to \R$}
        \begin{itemize}
            \item Inputs: $\vx \in \R^d$
            \item Params: weights $\vw \in \R^d$, bias $b \in \R$
            \item Output: $z = f(\vx) = \sum_i w_i x_i + b = \vw\T\vx + b$
        \end{itemize}
        \vskip 1em
\centering
        % Common/tikz/artificial_neuron.tex
\begin{tikzpicture}[scale=0.85]
\tikzstyle{every node}=[font=\scriptsize]
\tikzstyle{neuron} = [circle, draw, fill=thwsorange!20, minimum size=0.8cm, outer sep=3pt]
\tikzstyle{input}  = [circle, draw, fill=thwsblue!15, minimum size=0.5cm, outer sep=3pt]
\tikzstyle{conn}   = [thick, -Stealth, color=thwsblue]
\tikzstyle{outarr} = [thick, -Stealth]
\tikzstyle{annot} = [color=carnelian, font=\small]

% Input nodes
\node[input] (x1) at (-2.5,  0.9) {$x_1$};
\node[input] (x2) at (-2.5,  0.0) {$x_2$};
\node[input] (xd) at (-2.5, -1.4) {$x_d$};

% Neuron
\node[neuron] (n) at (0, 0) {$z$};
% Bias
\node[input, fill=thwsblue!5] (b1) at (0, 1.7) {$1$};

% Connections
\draw[conn] (x1) -- (n) node[midway, above] {$w_1$};
\draw[conn] (x2) -- (n) node[midway, above] {$w_2$};
\draw[conn] (xd) -- (n) node[midway, above] {$w_d$};
\draw[dotted, thick] (x2) -- (xd);
\draw[conn] (b1) -- (n) node[midway, right] {$b$};

% Output
\draw[outarr] (n) -- ++(1.0, 0.0) node[right] (z) {$z$};

% Labels
\node[annot] [below=0.1 of xd] {Inputs};
\node[annot] [below=0.1 of n] {Node};
\node[annot] [below=0.1 of z.south east] {Output};

\end{tikzpicture}
    \end{column}
\end{columns}
\end{frame}



% -------------------------------------------------------
% SECTION 3: From perceptron to MLP
% -------------------------------------------------------

% -------------------------------------------------------
% SECTION 4: Activation functions
% -------------------------------------------------------

% -------------------------------------------------------
% SECTION 5: Universal approximation
% -------------------------------------------------------

% -------------------------------------------------------
% SECTION 6: Notation summary
% -------------------------------------------------------

\end{document}